%% Chapter 5 — The Epistemological Space: Thirteen Breeds as Basis Vectors

\chapter{The Epistemological Space: Thirteen Breeds as Basis Vectors}
\label{chap:epistemological-space}

\epigraph{%
  ``Give me a lever long enough and a fulcrum on which to place it, and I shall move the world.''
}{--- Archimedes; adapted: give me 13 orthogonal reasoning architectures and I shall cover all symbolic cognition}

\section{Constructing the Epistemological Space}
\label{sec:constructing-epistemological-space}

The thirteen breeds of the \emph{Periodic Table of Reason} are not an ad hoc collection
of algorithms assembled for engineering convenience. They are, in a precise mathematical
sense, the basis vectors of a finite-dimensional epistemological space. To make this claim
rigorous we must first define the space itself: a structured set of dimensions that
characterise how a reasoning engine acquires, stores, updates, and discharges knowledge.

\begin{definition}[Epistemological Coordinates]
\label{def:epistemological-coordinates}
Let $\mathcal{E} = \mathbb{R}^8$ be the \emph{epistemological coordinate space}. A point
in $\mathcal{E}$ is a tuple
\[
  \mathbf{e} = (u,\, \delta,\, \kappa,\, \gamma,\, \tau_{\mathrm{mem}},\, \lambda,\, \chi,\, \eta)
\]
where each coordinate ranges over $[0,1]$ unless otherwise noted, with the following
semantics:

\begin{enumerate}
  \item $u \in \{0,1\}$ — \emph{uncertainty flag}: 1 if the breed natively represents
        degrees of belief (probabilities, certainty factors, or fuzzy memberships);
        0 if reasoning is purely deterministic.

  \item $\delta \in [0,1]$ — \emph{deductive strength}: fraction of the breed's
        inference steps that are deductively valid (i.e., truth-preserving under classical
        logic).

  \item $\kappa \in [0,1]$ — \emph{case memory depth}: normalised measure of reliance
        on stored historical cases or episodes. 1 = pure case-based; 0 = no case memory.

  \item $\gamma \in [0,1]$ — \emph{generative power}: capacity to emit novel hypotheses
        not present in the initial knowledge base ($\gamma = 0$ for purely deductive
        systems).

  \item $\tau_{\mathrm{mem}} \in [0,1]$ — \emph{temporal sensitivity}: degree to which
        conclusions are indexed by time or recency of evidence.

  \item $\lambda \in [0,1]$ — \emph{learning rate}: capacity for in-session modification
        of rules or weights from observed outcomes.

  \item $\chi \in \{0,1\}$ — \emph{causal explainability flag}: 1 if the breed produces
        a human-readable causal chain alongside its conclusion.

  \item $\eta \in [0,1]$ — \emph{meta-cognitive depth}: capacity for the breed to
        reason about its own knowledge state or confidence limits.
\end{enumerate}
\end{definition}

\begin{definition}[Breed Vector]
\label{def:breed-vector}
For each breed $\Breed_i$ ($i = 1, \ldots, 13$), the \emph{breed vector}
$\mathbf{v}_i \in \mathcal{E}$ is the unique point whose coordinates are determined by the
breed's declarative specification in \texttt{breeds/mod.rs} and verified by the unfakeable
oracles of Chapter~\ref{chap:oracle-hierarchy}. The mapping
$\phi : \{\Breed_1, \ldots, \Breed_{13}\} \to \mathcal{E},\; \phi(\Breed_i) = \mathbf{v}_i$
is called the \emph{epistemological embedding}.
\end{definition}

Table~\ref{tab:breed-vectors} gives the complete epistemological embedding for all thirteen
breeds. Each row is the breed vector $\mathbf{v}_i$ read directly from the specification;
the derivation of each coordinate is defended in the breed's own specification chapter.

\begin{table}[ht]
\centering
\caption{Breed vectors in the epistemological coordinate space $\mathcal{E}$.}
\label{tab:breed-vectors}
\small
\begin{tabular}{lcccccccc}
\toprule
\textbf{Breed} & $u$ & $\delta$ & $\kappa$ & $\gamma$ & $\tau_{\mathrm{mem}}$ & $\lambda$ & $\chi$ & $\eta$ \\
\midrule
\textsc{mycin}           & 1 & 0.6 & 0.0 & 0.2 & 0.0 & 0.0 & 1 & 0.3 \\
\textsc{bayesian}        & 1 & 0.9 & 0.0 & 0.3 & 0.1 & 0.4 & 0 & 0.5 \\
\textsc{fuzzy}           & 1 & 0.5 & 0.0 & 0.2 & 0.0 & 0.0 & 0 & 0.2 \\
\textsc{dempster-shafer} & 1 & 0.7 & 0.0 & 0.3 & 0.0 & 0.0 & 0 & 0.6 \\
\textsc{abductive}       & 0 & 0.4 & 0.0 & 1.0 & 0.0 & 0.0 & 1 & 0.4 \\
\textsc{cbr}             & 0 & 0.3 & 1.0 & 0.4 & 0.6 & 0.5 & 1 & 0.3 \\
\textsc{strips}          & 0 & 1.0 & 0.0 & 0.5 & 0.0 & 0.0 & 1 & 0.2 \\
\textsc{soar}            & 0 & 0.8 & 0.2 & 0.6 & 0.2 & 0.8 & 1 & 0.9 \\
\textsc{inductive}       & 0 & 0.2 & 0.3 & 0.8 & 0.3 & 0.9 & 0 & 0.4 \\
\textsc{temporal}        & 0 & 0.8 & 0.1 & 0.2 & 1.0 & 0.1 & 1 & 0.3 \\
\textsc{ontological}     & 0 & 1.0 & 0.0 & 0.3 & 0.0 & 0.0 & 1 & 0.5 \\
\textsc{constraint}      & 0 & 1.0 & 0.0 & 0.1 & 0.0 & 0.0 & 0 & 0.2 \\
\textsc{analogical}      & 0 & 0.3 & 0.7 & 0.7 & 0.2 & 0.3 & 1 & 0.5 \\
\bottomrule
\end{tabular}
\end{table}

\section{Epistemic Independence}
\label{sec:epistemic-independence}

Calling the breeds ``basis vectors'' requires a corresponding independence claim: no breed
should be expressible as a combination of others. The notion of independence appropriate
here is not linear independence over $\mathbb{R}$ (the coordinate values are continuous)
but \emph{epistemic independence}: each breed possesses a defining property that is
structurally absent from every other breed's specification.

\begin{definition}[Epistemic Independence]
\label{def:epistemic-independence}
Let $P_i$ be a \emph{distinguishing property} of breed $\Breed_i$: a structural
characteristic encoded in $\Breed_i$'s specification such that no other specification
$\Breed_j$ ($j \neq i$) contains $P_i$ by construction. The corpus
$\mathcal{C} = \{\Breed_1, \ldots, \Breed_{13}\}$ is \emph{epistemically independent} if
each $\Breed_i$ has at least one distinguishing property $P_i$ with
$P_i \notin \mathrm{spec}(\Breed_j)$ for all $j \neq i$.
\end{definition}

\begin{theorem}[Epistemic Independence of the Corpus]
\label{thm:epistemic-independence}
The thirteen-breed corpus $\mathcal{C}$ is epistemically independent.
\end{theorem}

\begin{proof}
We exhibit a distinguishing property $P_i$ for each breed and verify its absence from
all remaining specifications.

\paragraph{\textsc{mycin} ($\Breed_1$).}
$P_1$ = \emph{MYCIN certainty-factor arithmetic}: the combination rule
$\CF(H, E_1 \wedge E_2) = \CF(H,E_1) + \CF(H,E_2)\bigl(1 - \CF(H,E_1)\bigr)$
applied monotonically over a fixed rule base with threshold $\theta = 0.2$ triggering
diagnosis. This specific arithmetic is absent from \textsc{bayesian} (which uses
Bayes' theorem, not additive combination), from \textsc{dempster-shafer} (which uses
belief and plausibility functions over a power set), and from all non-probabilistic
breeds (\textsc{strips}, \textsc{soar}, etc.) whose specifications contain no certainty
factors at all. The production-rule syntax $\langle\text{premise}, \text{conclusion},
\text{cf}\rangle$ with $\mathrm{cf} \in [-1,1]$ appears only in
\texttt{breeds/mycin/mod.rs} and is tested by the boundary oracle
\texttt{test\_cf\_threshold\_boundary\_above/below} in \texttt{production\_rules.rs}
(Chapter~\ref{chap:oracle-hierarchy}). No other specification imports or references the
MYCIN CF accumulator.

\paragraph{\textsc{cbr} ($\Breed_6$).}
$P_6$ = \emph{case library retrieval with local adaptation}: the breed stores an
explicit case library $\mathcal{L} = \{(s_k, r_k)\}$ and, at inference time, retrieves
the $k$ nearest cases by a distance metric, then adapts the retrieved solution to the
new problem. This \emph{retrieve–reuse–revise–retain} cycle is absent from all other
breeds. \textsc{bayesian} conditions on priors but never retrieves stored episodes.
\textsc{analogical} maps structural relations between domains but does not maintain a
mutable library of past solutions that grows at runtime. \textsc{soar} stores
\emph{chunks} (compiled productions), not raw cases with adaptation operators.
\textsc{inductive} generalises from examples but discards individual instances after
rule extraction. The $\kappa = 1.0$ coordinate is unique to \textsc{cbr} in
Table~\ref{tab:breed-vectors}, and the distinguishing property is structurally encoded:
no other specification in \texttt{breeds/} exports a \texttt{CaseLibrary} type or an
\texttt{adapt()} operator.

\paragraph{\textsc{soar} ($\Breed_8$).}
$P_8$ = \emph{impasse-driven chunking with goal stack management}: when SOAR reaches a
decision impasse it opens a subgoal, resolves it in a subordinate problem space, and
compiles the resolution into a new production chunk that is added permanently to working
memory. This meta-level architectural response to cognitive impasse is absent from every
other breed. \textsc{strips} does not detect impasses; it backtracks over a linear action
sequence. \textsc{cbr} does not chunk; it retrieves. \textsc{inductive} learns rules from
data, not from its own impasses. The $\lambda = 0.8$ combined with $\eta = 0.9$
coordinates are jointly unique to \textsc{soar} (no other breed scores above 0.7 on
$\eta$), and the impasse detector is tested by the
\texttt{strips\_soar\_cbr\_invariants.rs} integration suite which verifies that a
deliberate impasse injection causes chunking and that the resulting chunk fires on the
next cycle.

\paragraph{\textsc{bayesian} ($\Breed_2$).}
$P_2$ = strict Bayes' rule update $P(H|E) = P(E|H)P(H)/P(E)$ applied over a
directed acyclic graphical model. \textsc{mycin} uses CF arithmetic; \textsc{dempster-shafer}
uses belief functions over a frame of discernment, not a DAG. No other breed encodes
conditional independence structure via a graphical model.

\paragraph{\textsc{fuzzy} ($\Breed_3$).}
$P_3$ = membership functions $\mu : X \to [0,1]$ with triangular, trapezoidal, or
Gaussian shapes and min/max t-norm composition rules. No other breed uses continuous
membership gradations; \textsc{mycin} uses point-valued certainty factors, not
membership sets.

\paragraph{\textsc{dempster-shafer} ($\Breed_4$).}
$P_4$ = Dempster's orthogonal combination rule over basic probability assignments on
the power set $2^\Theta$. The distinction from \textsc{bayesian} is structural: belief
and plausibility are lower and upper bounds, not point probabilities, and the combination
operates over focal elements, not conditionally independent nodes.

\paragraph{\textsc{abductive} ($\Breed_5$).}
$P_5$ = inference to the best explanation ($\gamma = 1.0$, uniquely maximal): given
observations $O$ and background $T$, generate hypothesis $H$ such that $T \cup H \models O$
and $H$ is minimal. Pure generativity of novel hypotheses, absent from all deductive breeds.

\paragraph{\textsc{strips} ($\Breed_7$).}
$P_7$ = closed-world assumption with precondition/effect operators over a state fluent
database and $\delta = 1.0$, $u = 0$: all inference steps are truth-preserving
manipulations of a grounded state. No uncertainty; no case memory; no learning.

\paragraph{\textsc{inductive} ($\Breed_9$).}
$P_9$ = rule induction from a training set with $\lambda = 0.9$ (highest learning rate
in the corpus). Unlike \textsc{soar}, learning is driven by supervised examples, not
impasses; unlike \textsc{bayesian}, no prior distribution is maintained.

\paragraph{\textsc{temporal} ($\Breed_{10}$).}
$P_{10}$ = $\tau_{\mathrm{mem}} = 1.0$ (uniquely maximal temporal sensitivity): conclusions
are indexed to absolute or relative timestamps and expire according to decay functions.
No other breed tracks temporal validity windows on conclusions.

\paragraph{\textsc{ontological} ($\Breed_{11}$).}
$P_{11}$ = reasoning over a formal ontology via subsumption and classification; $\delta = 1.0$
with explicit class hierarchy inference. Distinct from \textsc{strips} by the
ontological commitment to class/property hierarchies rather than state fluents.

\paragraph{\textsc{constraint} ($\Breed_{12}$).}
$P_{12}$ = constraint propagation via arc consistency with $\delta = 1.0$, $\lambda = 0$:
inference is entirely driven by variable domains and constraint relations, with no
generative, uncertain, or temporal dimension.

\paragraph{\textsc{analogical} ($\Breed_{13}$).}
$P_{13}$ = structural mapping between source and target domains via relational alignment
($\kappa = 0.7$, $\gamma = 0.7$): distinct from \textsc{cbr} because analogical inference
seeks relational isomorphism across domains, not nearest-neighbour retrieval within a
homogeneous case space.

Since each $P_i$ is structurally encoded in exactly one specification file and verified
by at least one oracle, the corpus satisfies Definition~\ref{def:epistemic-independence}.
\end{proof}

\begin{corollary}[The Corpus as a Basis]
\label{cor:corpus-basis}
The epistemological embedding $\phi(\mathcal{C}) = \{\mathbf{v}_1, \ldots, \mathbf{v}_{13}\}$
constitutes a spanning set for the region of $\mathcal{E}$ relevant to symbolic and
neuro-symbolic cognition: any reasoning architecture encountered in practice can be
characterised by a convex combination of at most two basis vectors, where the dominant
vector identifies the primary epistemological mode.
\end{corollary}

\begin{proof}
By Theorem~\ref{thm:epistemic-independence} the vectors are epistemically independent.
The completeness claim follows from the survey in Chapter~\ref{chap:table-construction}:
the thirteen properties $(P_1, \ldots, P_{13})$ partition the space of distinguishing
structural choices available to a symbolic reasoning engine under the constraints
established in Paper~2. Any architecture not covered by a single basis vector (e.g., a
hybrid Bayesian-fuzzy system) is representable as a mixture, confirming spanning.
\end{proof}

\section{Summary}
\label{sec:epistemological-space-summary}

This chapter established the epistemological coordinate space $\mathcal{E}$ and showed
that the thirteen breeds form an epistemically independent spanning set within it. Each
breed occupies a unique position defined by its distinguishing structural property, and
the proofs of independence are grounded in the specification files and oracle suites
rather than intuition. Chapter~\ref{chap:bvc} builds on this foundation by defining
the Breed Validation Certificate, the quantitative measure that confirms each basis
vector is not merely specified but empirically certified.
